\secnumbersection{MARCO CONCEPTUAL}

\subsection{LECTORES DE PANTALLA Y EL DOM}

Un lector de pantalla es un software de tecnología asistiva que convierte el contenido de una interfaz digital en voz sintetizada o en salida braille, permitiendo al usuario acceder a esa información sin necesidad de verla \cite{w3ctoolstechniques}. La salida braille se entrega mediante una línea braille refrescable: un dispositivo de hardware, conectado al computador, con celdas que forman y deshacen caracteres en relieve mediante pines accionados electromecánicamente, y que va mostrando al tacto y en tiempo real el contenido que el lector de pantalla comunica.

Los lectores de pantalla no son una categoría homogénea ni universal: existen embebidos en el propio sistema operativo, como Windows Narrator o VoiceOver en macOS/iOS, y otros que se instalan por separado, como JAWS o NVDA; dentro de estos últimos hay opciones de código abierto y sin costo (NVDA) y opciones pagas orientadas al uso profesional (JAWS); y hay tanto lectores de escritorio como lectores para dispositivos móviles, como TalkBack en Android. Qué lector de pantalla resulta más adecuado depende, entonces, del sistema operativo, el dispositivo y el presupuesto disponible, y no existe uno que cubra todos los contextos de uso por igual.

\begin{table}[h]
    \centering
    \caption{\label{table:lectorespantalla} Clasificación de los lectores de pantalla más usados.} Fuente: Elaboración propia, a partir de la Encuesta de Usuarios de Lectores de Pantalla \#10 de WebAIM \cite{webaim2024screenreadersurvey10}.
    \begin{tabular}{|p{2.4cm}|p{2.3cm}|p{2.8cm}|p{2.3cm}|p{4.2cm}|}
        \hline
        \textbf{Lector} & \textbf{Tipo} & \textbf{Licencia} & \textbf{Plataforma} & \textbf{Uso reportado} \\
        \hline
        JAWS & Instalado & Pago & Escritorio & 40,5\% como lector principal \\
        \hline
        NVDA & Instalado & Código abierto, sin costo & Escritorio & 37,7\% como lector principal \\
        \hline
        VoiceOver & Embebido (macOS/iOS) & Incluido con el sistema operativo & Escritorio y móvil & 9,7\% como lector principal de escritorio; 70,6\% de uso reportado en móviles \\
        \hline
        Narrator & Embebido (Windows) & Incluido con el sistema operativo & Escritorio & 0,7\% como lector principal \\
        \hline
        TalkBack & Embebido (Android) & Incluido con el sistema operativo & Móvil & 34,7\% de uso reportado en móviles \\
        \hline
    \end{tabular}
\end{table}

Las cifras de escritorio corresponden a la pregunta por el lector de pantalla \textbf{principal}, mientras que las cifras de uso en dispositivos móviles corresponden a una pregunta distinta por los lectores que el encuestado usa \textbf{con frecuencia}, sin exclusividad entre ellos \cite{webaim2024screenreadersurvey10}; por eso no son directamente comparables entre sí ni suman 100\%.

El DOM (\textit{Document Object Model}) es la representación estructurada del documento web: una jerarquía de nodos (elementos, atributos, texto) que el navegador construye a partir del código HTML de la página en su parte estática inicial, y sobre la que se puede consultar y operar mediante programación \cite{whatwgdom}. Esa segunda dimensión, programática, es la que aporta JavaScript: al ejecutarse, puede agregar, quitar o modificar nodos del DOM ya construido, de modo que lo que el navegador termina exponiendo no depende solo del HTML servido inicialmente, sino también de los cambios que JavaScript introduce después. El navegador deriva del DOM resultante un árbol de accesibilidad, una representación paralela y más acotada que expone únicamente los roles, nombres, estados y relaciones relevantes para una tecnología asistiva, a través de la API de accesibilidad del sistema operativo (por ejemplo, UI Automation en Windows o AXAPI en macOS). Es esa API la que el lector de pantalla consulta para construir lo que finalmente comunica al usuario.

La capacidad del navegador de construir un árbol de accesibilidad completo depende de que la página declare correctamente su semántica, algo que el HTML sin marcadores adicionales no siempre garantiza. El W3C define para ello dos estándares complementarios: las Pautas de Accesibilidad para el Contenido Web (WCAG), que establecen los criterios que debe cumplir un sitio para considerarse accesible \cite{wcag22}, y WAI-ARIA, un conjunto de atributos (roles, estados y propiedades \texttt{aria-*}) que permite declarar la semántica de elementos que no la expresan por sí solos —por ejemplo, marcar una notificación emergente como una región viva (\texttt{aria-live}) para que se anuncie apenas aparece— \cite{waiaria2023}. El atributo \texttt{alt}, por su parte, es el mecanismo nativo para asociar un texto alternativo a una imagen —y existe desde antes que los lectores de pantalla, no exclusivamente para ellos—; sin él, una imagen no tiene ninguna representación textual que el lector de pantalla pueda comunicar \cite{w3caltimages}.

El lector de pantalla permite al usuario recorrer \textbf{secuencialmente} el contenido de la página mediante la tecla Tab, uno tras otro: no solo los elementos interactuables como campos de texto o enlaces, sino también contenido no interactuable como párrafos de texto. Ese orden de recorrido no es estrictamente el orden en que los elementos fueron declarados en el HTML: el atributo \texttt{tabindex} permite alterarlo, adelantando o postergando elementos respecto de su posición original en el documento. Además de ese recorrido, el lector de pantalla ofrece ayudas de navegación que no existen sin él: una de las más usadas es el listado de enlaces, que presenta todos los enlaces de la página en una lista y permite saltar directamente al que se necesite.

La eficacia de este recorrido secuencial depende de la familiaridad previa del usuario con la estructura de la página: sobre un sitio ya explorado, permite ubicar rápidamente los elementos buscados; frente a un sitio nuevo, requiere recorrerlo y estudiarlo varias veces antes de poder anticipar su estructura.

\subsection{EXTENSIONES DE NAVEGADOR}

Una extensión es un módulo de software que un usuario instala sobre su navegador para modificar o ampliar su comportamiento —por ejemplo, para bloquear anuncios, gestionar contraseñas o capturar información de una página—, sin ser una aplicación independiente ni una función nativa del propio navegador \cite{mdnwebextensions}.

Google Chrome fue el primer navegador en incorporar un sistema de extensiones basado enteramente en HTML, CSS y JavaScript, habilitado en su canal de desarrollo en 2009 y lanzado de forma pública junto a Chrome 4.0 en enero de 2010 \cite{chromeextensionshistory}. A partir de esa base, Mozilla adoptó en 2015 una API compatible para su navegador Mozilla Firefox, bautizada WebExtensions, que buscaba convertirse en un estándar común entre navegadores; en 2021 se llegó a formar un grupo comunitario del W3C específicamente para intentar formalizar esa base común entre distintos fabricantes \cite{w3cwebextensionscg2021}, pero ese intento de estandarización no llegó a concretarse. Pese a ello, la popularidad de Chrome llevó a que el resto de los navegadores basados en su motor Chromium (Edge, Opera, Brave, entre otros) convergieran hacia una API de extensiones muy similar \cite{mdnwebextensions}.

Hoy existen, en términos generales, dos grandes familias de navegador incompatibles entre sí: los basados en Chromium, que agrupan a la gran mayoría, y Mozilla Firefox, basado en el motor Gecko y heredero del antiguo navegador Netscape. Aparte de estas dos familias existen otros navegadores, como los propietarios de ciertos fabricantes de dispositivos móviles. Un caso aparte es Safari que utiliza su propio motor, WebKit. Independiente de que un navegador soporte extensiones, estas solo son compatibles con el navegador para el que fueron diseñadas: el código y el empaquetado de una extensión para Chrome no son directamente instalables en Firefox ni viceversa, aún cuando ambas compartan la mayoría de las llamadas de su API.

A nivel de capacidades, una extensión se ejecuta bajo un modelo de permisos explícitos declarados en un archivo de manifiesto: solo puede acceder a las páginas, pestañas o funciones del navegador para las que fue autorizada, y se organiza mediante scripts de contenido —que se ejecutan en el contexto de la página visitada— y un script de fondo, que coordina la lógica de la extensión y sus comunicaciones con servicios externos. Ese modelo de permisos delimita tanto lo que la extensión puede automatizar dentro del navegador como lo que queda fuera de su alcance sin intervención directa del usuario o del sistema operativo.

\subsection{RECONOCIMIENTO Y SÍNTESIS DE VOZ (STT Y TTS)}

El reconocimiento automático del habla (\textit{automatic speech recognition}, ASR; también conocido como \textit{speech to text}, STT) es el proceso mediante el cual un sistema decodifica una señal de audio con voz humana y la transcribe a texto, típicamente combinando un modelo acústico —que traduce la señal a una secuencia probable de fonemas— y un modelo de lenguaje —que determina, a partir de esos fonemas, las palabras más probables— \cite{ibmasr}.

La síntesis de voz (\textit{text to speech}, TTS) es el proceso inverso: toma una entrada textual y la convierte en audio hablado. Esa entrada no siempre es texto plano: estándares como SSML (\textit{Speech Synthesis Markup Language}), del W3C, permiten anotar el texto con marcadores de énfasis, acento, velocidad o pausas, para un control más fino de cómo se pronuncia \cite{w3cssml}. A nivel de técnica, la síntesis de voz se ha resuelto históricamente de más de una forma: por concatenación de fragmentos de voz pregrabada, por síntesis paramétrica a partir de un modelo estadístico de la voz, y, más recientemente, mediante modelos neuronales generativos que producen la forma de onda directamente a partir de datos de entrenamiento \cite{tan2021neuraltts}.

Lectores de pantalla como JAWS o NVDA no necesariamente implementan su propio motor de síntesis: en general, delegan la conversión final de texto a voz a un motor TTS del sistema operativo o a una voz de terceros instalada por el usuario, y lo que le corresponde al lector de pantalla es determinar qué texto de la interfaz enviar a ese motor.

\subsection{IA MULTIMODAL: MODELOS DE LENGUAJE VISUAL (VLM)}

\subsubsection{IA generativa de lenguaje}

Los modelos de lenguaje generativos son un tipo de inteligencia artificial capaz de producir texto nuevo a partir de patrones aprendidos de grandes volúmenes de datos, en lugar de limitarse a clasificar o predecir una etiqueta sobre texto ya existente \cite{ibmgenpredictive}. Esa generación se implementa técnicamente como una predicción: un modelo autoregresivo —que genera cada nueva unidad de texto a partir de lo que él mismo generó previamente—, como los que se basan en la arquitectura Transformer —una arquitectura de red neuronal que pondera mediante mecanismos de atención qué partes de la entrada son más relevantes para cada predicción—, no produce una respuesta completa de una sola vez: predice, \textit{token} a \textit{token} —el \textit{token} es la unidad mínima de texto en que el modelo divide el lenguaje, por ejemplo una palabra o parte de ella—, cuál es la unidad de texto más probable dado todo lo que ha recibido y generado hasta ese momento, y repite ese proceso hasta completar la respuesta \cite{vaswani2017attention}. Es decir, ``generar'' texto y ``predecir la palabra siguiente'' no son dos capacidades distintas, sino una relación de mecanismo a resultado: el modelo predice de forma iterativa, \textit{token} a \textit{token}, y esa cadena de predicciones es percibida por quien lo usa como la generación de una respuesta nueva.

\subsubsection{IA multimodal}

La IA multimodal corresponde a modelos de IA capaces de procesar y/o producir información de más de un tipo de dato —texto, imágenes, audio o video— de forma conjunta, en lugar de limitarse a un único tipo de entrada o salida \cite{wu2023multimodalsurvey}. Dentro de esa familia se encuentran los modelos de lenguaje visual (VLM, \textit{vision language model}), especializados en la interacción entre texto e imágenes: reciben como entrada una combinación de imagen(es) y texto, y devuelven una salida en texto, apoyándose en una representación interna que alinea los conceptos visuales de la imagen con los conceptos del lenguaje \cite{wu2023multimodalsurvey}.

\subsubsection{Entrenamiento}

El entrenamiento es el proceso mediante el cual un modelo de IA aprende a reconocer patrones y correlaciones para optimizar su desempeño en tareas similares a las de su uso previsto \cite{ibmmodeltraining}. Ese aprendizaje se realiza sobre un \textit{dataset}: una colección de datos de ejemplo relevantes para la tarea que el modelo debe resolver —por ejemplo, pares de pregunta y respuesta, o imágenes junto a su descripción—; cuanto más se asemeje ese conjunto a los casos reales que el modelo enfrentará después, mejor será su desempeño sobre datos nuevos \cite{ibmtrainingdata}.

Un entrenamiento puede fallar en dos direcciones opuestas. Por sobreajuste (\textit{overfitting}), el modelo se ajusta tan estrechamente a los ejemplos de entrenamiento que termina memorizándolos en lugar de generalizar, logrando buenos resultados solo sobre datos ya vistos; por subajuste (\textit{underfitting}), el modelo es demasiado simple, o el entrenamiento demasiado breve, para siquiera capturar los patrones del propio conjunto de entrenamiento \cite{ibmoverfitting}. Un entrenamiento adecuado, en cambio, encuentra un punto intermedio entre ambos extremos, que le permite generalizar razonablemente bien a datos nuevos de una naturaleza similar a los ya vistos.

De ahí se desprende una limitación central de cualquier modelo de IA: no puede responder de forma confiable sobre información que su \textit{dataset} de entrenamiento no consideró, o que consideró en una cantidad insuficiente para generalizar ese patrón sin sesgo. Cuando un tipo de caso está subrepresentado en el \textit{dataset}, el modelo tiende a fallar o a generalizar incorrectamente para ese caso a partir de lo aprendido sobre los casos mayoritarios \cite{ibmdatabias}. En términos simples, un modelo entrega mejores resultados sobre aquello que aparece con frecuencia en su \textit{dataset} de entrenamiento, y resultados progresivamente peores sobre aquello que aparece con poca frecuencia o no aparece en él.

\subsubsection{Inferencia}

La inferencia es el acto de usar un modelo de IA ya entrenado para obtener una salida a partir de una entrada nueva \cite{nvidiaaiinference}. Ejecutar esa inferencia sobre un modelo de gran tamaño es, en sí mismo, un desafío computacional: un equipo de escritorio promedio no puede hacerlo, o al menos no en tiempos que permitan un uso práctico, por lo que habitualmente se recurre a hardware específico —típicamente GPU— para hacerla eficiente. Frente a esto existen dos caminos habituales, ambos con inconvenientes: los servicios de IA en línea corren sobre datacenters de gran escala y cobran por consulta; además, procesan los datos del usuario en infraestructura de terceros, sin garantías claras sobre su tratamiento.

En este contexto aparece el concepto de latencia: el tiempo de espera entre que se envía una consulta y se recibe la respuesta. Es un factor crítico para la experiencia de uso, ya que una espera prolongada puede interrumpir el flujo de la interacción del usuario, incluso cuando el contenido de la respuesta es correcto.

\subsection{EVALUACIÓN AUTOMATIZADA DE VLMs}

Existe una gran oferta de modelos VLM \textit{open source}, y distintos modelos entregan resultados diferentes para las mismas consultas, por lo que elegir el más pertinente para una necesidad determinada no es trivial. Hacerlo de forma manual —tomar un VLM, ingresarle N preguntas y evaluar cada una de esas N respuestas— no escala: para que el resultado sea confiable se necesita un N grande, y probar varios modelos vuelve la evaluación manual una tarea repetitiva e inviable. De ahí la necesidad, documentada en la literatura, de automatizar este proceso mediante técnicas de \textit{benchmarking} sistemático.

\subsubsection{Rúbrica cualitativa}

Una rúbrica es un instrumento de evaluación cualitativa que describe, para cada nivel de desempeño posible, los criterios concretos que una respuesta debe cumplir para alcanzarlo; se construye de forma manual antes de la evaluación y permite calificar de forma consistente incluso cuando la respuesta es abierta y no tiene una única forma correcta —por ejemplo, al calificar la respuesta de un estudiante en un examen de argumentación, donde una rúbrica puede distinguir niveles que van desde una postura sin fundamento hasta un argumento completo y bien sustentado. La evaluación de un VLM enfrenta el mismo problema de fondo: la respuesta a una consulta sobre una imagen tampoco tiene una única forma correcta, por lo que una rúbrica permite distinguir niveles de calidad —por ejemplo, si identifica los elementos relevantes de la imagen, si omite información importante, o si incluye contenido que no está presente en ella— en lugar de exigir una coincidencia exacta con una respuesta de referencia.

\subsubsection{\textit{LLM-as-a-judge}}

``\textit{LLM-as-a-judge}'' es un paradigma de evaluación en el que un LLM, en lugar de un evaluador humano, cumple el rol de juez: se le entrega la respuesta que se quiere evaluar junto con los criterios de evaluación —por ejemplo, una rúbrica—, y el modelo determina si esos criterios están presentes o ausentes para asignar un puntaje de forma cuantitativa \cite{zheng2023llmasjudge}.

\subsubsection{\textit{LLM-as-a-judge-evaluation}}

Para confirmar que un LLM evaluador funciona bien, en este trabajo se introduce el concepto de ``\textit{LLM-as-a-judge-evaluation}'': un humano califica de forma independiente las mismas respuestas que evalúa el LLM, y luego se contrasta cuánto se acercan los puntajes de uno y otro. Es una segunda capa de evaluación, destinada a comprobar que el criterio con el que el LLM juzga esas respuestas está bien calibrado.

Con la \textit{LLM-as-a-judge-evaluation} se cierra el ciclo: el humano evalúa las respuestas, el LLM las evalúa también, y la comparación entre ambos permite validar al LLM evaluador a partir del criterio humano.

\subsubsection{Herramientas estandarizadas de evaluación}

Respecto de los sistemas estandarizados de evaluación específica para VLMs, ya existen herramientas y métricas destinadas a esta tarea. Un ejemplo es VLMEvalKit (open-compass): un \textit{toolkit open source} de evaluación de modelos multimodales que unifica la evaluación de una gran cantidad de modelos sobre decenas de \textit{benchmarks} con configuración mínima, automatizando la preparación de datos, la inferencia y el cálculo de métricas, e incluso usando un LLM como extractor de respuestas cuando la comparación exacta falla \cite{duan2024vlmevalkit}. Sin embargo, ninguna de estas herramientas está diseñada específicamente para la lectura de sitios web en el contexto de la discapacidad visual, que es el foco de este trabajo.